\scanpage{48}
\mspage{3}
On a \unsure{les|des} équivalences de catégories
\diagram{
\begin{tikzcd}[ampersand replacement=\&]
(13)\quad \operatorname{Rev\,\text{ét}}(\widetilde\Sigma) \arrow[r,"\approx","\text{équiv.}"'] \& \operatorname{Rev\,\text{et}}(\widetilde\Sigma^\circ) \arrow[r,"\sim","\text{iso}"'] \& \operatorname{Rev\,\text{et}}(\Sigma^*)
\end{tikzcd}
}
déduites de l'inclusion
\[
\widetilde\Sigma^\circ \overset{\mathrm{def}}= \widetilde\Sigma\setminus\unsure{\partial|\delta}\widetilde\Sigma = \widetilde\Sigma|\Sigma^* \hookrightarrow \widetilde\Sigma
\]
et de l'isomorphisme
\[
\widetilde\Sigma^\circ \xrightarrow{\sim} \Sigma^*,
\]
d'où une équivalence entre les groupoïdes fondamentaux au sens large (\unsure{dt|dont} les objets sont les rev. étales 1-connexes)\deleted{\unsure{des}}
\[
(14)\qquad \Pi_1(\widetilde\Sigma) \longrightarrow \Pi_1(\Sigma^*)
\]
et en composant avec
\[
\Pi_1(U_i)\longrightarrow \Pi_1(\widetilde\Sigma)
\]
déduit de l'inclusion $U_i\hookrightarrow\widetilde\Sigma$, on trouve
\[
(15)\qquad \Pi_1(\deleted{\unsure{\Sigma}}U_i)\longrightarrow\Pi_1(\Sigma^*).
\]
Nous allons représenter graphiquement l'homomorphisme 14, en ``dessinant'' $\widetilde\Sigma$ sur $\Sigma^*$, comme une région qui serait contenue dans $\Sigma^*$ en lui enlevant \inserted{l'intérieur de} trois disques \inserted{fictifs} $\unsure{\Delta_i}$ ``infiniment petits'' \deleted{et} (s'identifiant aux disques ouverts des $T_{\Sigma,R_i}$ \deleted{\unsure{g}}, ou encore aux cônes sur les $U_i$), et on se représente (14) comme provenant d'une telle inclusion fictive.
\margin{Plus précisément au lieu de travailler sur $\Sigma^*$, nous travaillons sur $C$ \unsure{``sphère à trous''} un peu plus grande que $\widetilde\Sigma$.}
\margin{47}
